\chapter{Bacteria and Viruses}\label{ch:g8:microbiology}

\omterm{def:g4:microbes:def}{Microbes} are everywhere: soil, \omterm{def:g4:plants-need:needs}{water}, your \omterm{def:g5:nerves:senses}{skin} and gut. Most are harmless
or helpful; some cause disease. This chapter contrasts \omterm{def:g4:microbes:def}{bacteria} and
viruses and introduces safe thinking about infection.

\section{Bacteria}

\begin{definition}[Bacterium]\label{def:g8:microbiology:bacteria}
A \emph{bacterium}\index{bacterium} (plural \omterm{def:g4:microbes:def}{bacteria}) is a prokaryotic
microorganism: no \omterm{def:g8:organelles:nucleus}{nucleus}, DNA in a nucleoid, \omterm{def:g8:organelles:wall}{cell wall} (often), and
\omterm{def:g5:reproduction:def}{reproduction} mainly by binary fission. Sizes are typically around
\qty{1}{\micro m}.
\end{definition}
\omimg{microscope.jpg}{0.50\linewidth}{Microbes are often studied with a microscope.}


\begin{definition}[Useful roles]\label{def:g8:microbiology:useful}
\omterm{def:g4:microbes:def}{Bacteria} decompose dead matter, fix nitrogen in some \omterm{def:g1:plants-animals:plant}{plant} partnerships,
help digest \omterm{def:g2:food-energy:food}{food} in guts, and are used in \omterm{def:g2:food-energy:food}{food} production (yoghurt,
cheese) and biotechnology.
\end{definition}

\begin{omfigure}
\begin{tikzpicture}[font=\footnotesize]
  \draw[thick, omDef, fill=omDef!15] (0,0) ellipse (1.3 and 0.55);
  \fill[omThm!50] (-0.2,0) circle (0.25);
  \node at (0,-1.0) {rod (bacillus)};
  \draw[thick, omDef, fill=omDef!15] (3.5,0) circle (0.7);
  \fill[omThm!50] (3.5,0) circle (0.2);
  \node at (3.5,-1.0) {sphere (coccus)};
  \draw[thick, omDef, fill=omDef!15] (7.0,0) -- (7.5,0.4) -- (8.0,0)
    -- (7.5,-0.4) -- cycle;
  \draw[thick, omDef, fill=omDef!15] (7.3,0.15) -- (7.8,0.55) -- (8.3,0.15)
    -- (7.8,-0.25) -- cycle;
  \node at (7.6,-1.0) {spiral forms};
\end{tikzpicture}

{\small Common bacterial shapes (schematic). Shape alone does not tell you
if a \omterm{def:g6:classification:species}{species} is harmful.}
\end{omfigure}

\section{Viruses}

\begin{definition}[Virus]\label{def:g8:microbiology:virus}
A \emph{virus}\index{virus} is an infectious particle with genetic
material (DNA or RNA) inside a \omterm{def:g6:nutrition:groups}{protein} coat (sometimes with a membrane
envelope). \omterm{def:g4:microbes:def}{Viruses} are not \omterm{def:g7:microscope:cell}{cells}: they cannot \omterm{def:g2:what-alive:signs}{reproduce} or metabolise
independently --- they hijack host \omterm{def:g7:microscope:cell}{cells}.
\end{definition}

\begin{omfigure}
\begin{tikzpicture}[font=\footnotesize]
  \draw[thick, omProp, fill=omProp!20] (0,0) circle (0.9);
  \fill[omThm!40] (0,0) circle (0.35);
  \node[align=center] at (0,-1.4) {simple virus\\(genome + coat)};
  \draw[thick, omExo, fill=omExo!15] (4.5,0) -- (4.5,1.2) -- (5.2,1.6)
    -- (5.9,1.2) -- (5.9,0) -- cycle;
  \draw[thick, omExo] (5.2,-0.1) -- (5.2,-0.9);
  \draw[thick, omExo] (4.7,-0.9) -- (5.7,-0.9);
  \node[align=center] at (5.2,-1.5) {bacteriophage\\(virus of bacteria)};
\end{tikzpicture}

{\small \omterm{def:g4:microbes:def}{Viruses} come in many architectures; all need a host to replicate.}
\end{omfigure}

\begin{proposition}[Bacteria vs viruses]\label{prop:g8:microbiology:compare}
\begin{itemize}
  \item \omterm{def:g4:microbes:def}{Bacteria}: \omterm{def:g7:microscope:cell}{cells}; can \omterm{def:g2:what-alive:signs}{grow} on suitable media; many treatable with
        \emph{antibiotics}\index{antibiotic} (which target bacterial
        processes).
  \item \omterm{def:g4:microbes:def}{Viruses}: non-cellular; need living hosts; \omterm{prop:g8:microbiology:compare}{antibiotics} do not kill
        viruses; prevention often relies on hygiene and
        \emph{vaccines}\index{vaccine}.
\end{itemize}
\end{proposition}
\admitted

\section{Infection and defence (preview)}

\begin{definition}[Pathogen]\label{def:g8:microbiology:pathogen}
A \emph{pathogen}\index{pathogen} is a disease-causing \omterm{def:g4:ecosystems:levels}{organism} or
infectious agent (certain \omterm{def:g4:microbes:def}{bacteria}, viruses, \omterm{def:g4:microbes:def}{fungi}, parasites).
\end{definition}

\begin{example}[Transmission]\label{ex:g8:microbiology:trans}
\omterm{def:g8:microbiology:pathogen}{Pathogens} spread by droplets, \omterm{def:g1:senses:five}{touch}, contaminated \omterm{def:g2:food-energy:food}{food}/\omterm{def:g4:plants-need:needs}{water}, vectors
(e.g.\ mosquitoes) or \omterm{def:g7:digestion:blood}{blood}. Breaking transmission (handwashing, clean
\omterm{def:g4:plants-need:needs}{water}, safe \omterm{def:g2:food-energy:food}{food}) is public \omterm{def:g2:healthy:health}{health} in action.
\end{example}

\begin{method}[Safe lab habits with microbes]\label{met:g8:microbiology:safe}
\begin{enumerate}
  \item Treat all cultures as potentially hazardous.
  \item Aseptic technique: sterilise loops, limit dish opening time.
  \item Never culture unknown environmental samples from risky \omterm{def:g1:caring:needs}{places}
        without specialist rules.
  \item Disinfect benches; wash \omterm{def:g1:my-body:parts}{hands}; dispose of waste as instructed.
\end{enumerate}
\end{method}

\section{Exercises}

\begin{exercise}[$\star$]\label{exo:g8:microbiology:1}
Are \omterm{def:g4:microbes:def}{bacteria} prokaryotes or eukaryotes? What does that mean for the
\omterm{def:g8:organelles:nucleus}{nucleus}?
\end{exercise}

\begin{exercise}[$\star$]\label{exo:g8:microbiology:2}
Why are viruses not considered \omterm{def:g7:microscope:cell}{cells}?
\end{exercise}

\begin{exercise}[$\star$]\label{exo:g8:microbiology:3}
Give two useful roles of \omterm{def:g4:microbes:def}{bacteria}.
\end{exercise}

\begin{exercise}[$\star$]\label{exo:g8:microbiology:4}
Why do \omterm{prop:g8:microbiology:compare}{antibiotics} not cure a typical viral cold?
\end{exercise}

\begin{exercise}[$\star$]\label{exo:g8:microbiology:5}
Define \omterm{def:g8:microbiology:pathogen}{pathogen} and give one bacterial and one viral example (names at
everyday level are fine).
\end{exercise}

\section{Problem: Outbreak thinking}

